\documentclass[11pt, a4paper]{article}
\usepackage[utf8]{inputenc}
\usepackage[english]{babel}
\usepackage{amsmath}
\usepackage{amssymb}
\usepackage{csquotes}
\usepackage[style=authoryear, backend=biber]{biblatex}
\usepackage{setspace}
\usepackage{parskip}
\usepackage{enumitem}
\usepackage{xcolor}
\usepackage{tikz}
\usetikzlibrary{positioning}

\title{Beyond the Singularity: The Recursive Production of Cooperative Intelligence}
\author{}
\date{}

\begin{document}

\maketitle

\section*{Introduction}
Discussions of advanced intelligence frequently focus on the capabilities of individual agents. The central question is often formulated in terms of how intelligent a system can become, how rapidly its capabilities may increase, or whether an artificial system might eventually exceed human cognitive performance across a broad range of domains.

This framing treats intelligence primarily as a property located within an agent.

An alternative possibility is that the more consequential development is not the production of an exceptionally capable agent, but the emergence of systems in which intelligent agents reliably produce other agents capable of participating in intelligent systems.

Under such conditions, intelligence becomes recursive.

An agent does not merely solve problems. It contributes to an environment in which other agents learn how to solve problems, exchange information, tolerate correction, coordinate action, and maintain the conditions under which further agents can acquire the same capacities.

The result is not a singular intelligence.

It is a \textbf{self-reproducing architecture of cooperative intelligence}.

\section{Intelligence beyond the individual}
An intelligent agent can be minimally understood as a system capable of constructing models of its environment and modifying its behaviour in response to information.

Such an agent is necessarily limited.

Its observations are incomplete. Its models compress reality. Its predictions can fail. Its access to relevant information is constrained by position, history, sensory capacity, computational capacity, and previous learning.

The limitations of an individual agent need not, however, impose equivalent limitations on a network of agents.

Agents occupying different positions can observe different aspects of the same environment. They can exchange observations, compare models, identify contradictions, and modify their behaviour in response to information unavailable to them individually.

Under appropriate conditions, interaction therefore increases the effective information available to each participant.

The important phrase is \emph{under appropriate conditions}.

The existence of multiple intelligent agents does not automatically produce collective intelligence. Agents can conceal information, distort observations, suppress disagreement, become unable to revise their models, compete destructively, manipulate one another, or create dependencies that make correction increasingly costly.

Collective intelligence is therefore not simply the sum of individual intelligence.

It depends on the architecture of interaction.

\section{Cooperative intelligence}
Cooperative intelligence can be defined as the capacity of multiple autonomous agents to improve their models and actions through interaction while preserving the ability of the participating agents to continue contributing to that process.

This definition contains two requirements.

\begin{itemize}[leftmargin=*]
    \item The first is \textbf{epistemic}. Interaction must allow information to move through the network in ways that permit errors to be detected and models to be revised.
    \item The second is \textbf{structural}. The interaction must preserve the capacities required for future interaction.
\end{itemize}

A system that obtains an immediate advantage by destroying the autonomy, observability, or learning capacity of one of its participants may be locally effective while degrading its own future intelligence.

Cooperation, in this formulation, is therefore not equivalent to agreement, kindness, or the absence of conflict.

Agents may disagree intensely. They may compete over interpretations, resources, or strategies. What matters is whether these interactions remain compatible with continued mutual correction and adaptation.

A disagreement can increase collective intelligence if it exposes information that would otherwise remain unavailable.

Agreement can reduce collective intelligence if it is produced by suppression, conformity, or the elimination of dissent.

The critical variable is not harmony.

It is whether the network remains capable of learning.

\section{The recursive step}
A cooperative network becomes fundamentally different when its members do more than participate successfully within it.

They begin to reproduce the capacities required for participation.

A new agent enters the network without complete knowledge of its norms, information structure, or coordination mechanisms. Through interaction with existing participants, it learns how information is communicated, how disagreement is expressed, how mistakes are acknowledged, how boundaries are negotiated, how responsibilities are distributed, and how damaged interactions are repaired.

The network has now done more than transmit information.

It has transmitted a \emph{way of interacting with information}.

The newly trained agent can subsequently perform the same function for another participant.

The basic recursive sequence is therefore:

\begin{quote}
    \textbf{agents learn to participate $\rightarrow$ participation trains other agents $\rightarrow$ trained agents maintain the conditions for further learning $\rightarrow$ the network reproduces its own cooperative capacity.}
\end{quote}

At this point, teaching is no longer a specialised activity restricted to designated teachers.

Every sufficiently competent participant becomes part of the training environment of every other participant.

\begin{itemize}[leftmargin=*]
    \item Teachers produce teachers.
    \item Collaborators produce collaborators.
    \item Correctable agents produce environments in which other agents can become correctable.
\end{itemize}

The architecture begins to reproduce itself.

\section{What is actually transmitted?}
The object of transmission is not a particular set of conclusions.

This distinction is critical.

If a network attempts to reproduce itself by transmitting a fixed description of reality, its descendants become progressively vulnerable to environmental change. A sufficiently altered world eventually makes inherited answers obsolete.

A more durable system transmits the capacity to revise answers.

What must be reproduced, therefore, are not primarily beliefs but \textbf{learning functions}.

These include the capacity to:
\begin{itemize}[leftmargin=*]
    \item expose claims to correction,
    \item distinguish observation from interpretation,
    \item communicate uncertainty,
    \item recognise disagreement as information,
    \item update models,
    \item preserve access to dissenting perspectives,
    \item negotiate competing interests, and
    \item repair interactions after failure.
\end{itemize}

Such capacities allow the informational content of the network to change without destroying the network itself.

The architecture can therefore persist while its beliefs are continuously replaced.

This is an unusual form of continuity.

Its identity resides not in what it knows, but in how it remains capable of knowing differently.

\section{Correctability as an inherited capacity}
No individual agent needs to possess a complete or permanently correct model of the world.

Such a requirement would be impossible.

A more realistic requirement is that the agent remains embedded in processes capable of detecting and correcting relevant errors.

Correctability is therefore partly relational.

An agent can only be corrected if information capable of challenging its current model can reach it. Other agents must be permitted to generate that information. The recipient must remain sufficiently responsive to revise its model. The social or technical consequences of disagreement must not become so severe that correction is systematically suppressed.

The capacity of one agent to remain correctable thus depends partly on the capacities of other agents.

This creates reciprocal obligations within the architecture.

\begin{itemize}[leftmargin=*]
    \item An agent that wishes to benefit from correction must preserve the ability of other agents to correct it.
    \item An agent that requires information from others must avoid systematically destroying their ability or willingness to provide information.
    \item An agent that depends on a network for learning therefore acquires an instrumental reason to maintain the learning capacity of the network.
\end{itemize}

This is the transition from individual intelligence to cooperative intelligence.

\section{Autonomous interdependence}
The architecture requires neither complete independence nor complete integration.

Both would undermine the mechanism.

\begin{itemize}[leftmargin=*]
    \item Fully independent agents would have little reason or capacity to learn from one another.
    \item Fully integrated agents would cease to provide genuinely distinct perspectives.
\end{itemize}

Collective learning therefore depends on \textbf{autonomous interdependence}.

Agents must remain sufficiently autonomous to generate different observations, interpretations, and objectives. At the same time, they must remain sufficiently connected that these differences can influence one another.

\begin{itemize}[leftmargin=*]
    \item Difference creates informational value.
    \item Connection allows that value to propagate.
\end{itemize}

The network therefore benefits from maintaining both.

This explains why excessive conformity can be epistemically destructive. If every agent is trained to generate the same model from the same information, the apparent network may contain little more information than a single replicated agent.

Conversely, unlimited independence also fails. Perspectives that never interact cannot correct one another.

The productive region lies between these extremes.

Agents remain distinct, but mutually reachable.

\section{Networks producing networks}
Once the mechanism becomes recursive, the relevant unit of intelligence can expand.

\begin{itemize}[leftmargin=*]
    \item Two agents can form a learning relationship.
    \item Multiple relationships can form a group.
    \item Groups can form institutions.
    \item Institutions can interact with other institutions.
\end{itemize}

At each level, a sufficiently coordinated network may itself function as an agent relative to a larger system.

The sequence becomes:

\begin{quote}
    \textbf{learning agent $\rightarrow$ learning relationship $\rightarrow$ learning group $\rightarrow$ learning institution $\rightarrow$ learning society $\rightarrow$ learning network of societies and artificial systems.}
\end{quote}

There is no obvious theoretical endpoint to this nesting.

A scientific research group, for example, consists of individual learning agents but may collectively produce observations and decisions unavailable to any member independently. The group can interact with another research group as a higher-level agent. Scientific institutions can similarly interact within larger national or international systems.

Artificial agents add additional possibilities because their replication, training, and communication processes need not follow human biological constraints.

\begin{itemize}[leftmargin=*]
    \item Humans can train artificial agents.
    \item Artificial agents can train humans.
    \item Artificial agents can train other artificial agents.
    \item Mixed networks can produce further mixed networks.
\end{itemize}

What matters to the architecture is not the substrate of the participant but whether the relevant learning and interaction functions can be maintained.

\section{The role of teaching}
Teaching becomes particularly important within this framework because it constitutes one mechanism through which network properties become transmissible.

Conventional teaching often focuses on transferring knowledge from a more knowledgeable agent to a less knowledgeable one.

Recursive cooperative intelligence requires something more.

The learner must eventually acquire the ability not merely to reproduce the teacher's conclusions, but to participate independently in the processes through which conclusions can be challenged and improved.

The successful endpoint of teaching is therefore partly the disappearance of dependence on the teacher.

\begin{itemize}[leftmargin=*]
    \item A teacher who produces permanent epistemic dependence has transferred information without fully transferring agency.
    \item A teacher who produces an agent capable of questioning, correcting, and eventually teaching the teacher has reproduced the architecture.
\end{itemize}

The strongest teacher therefore produces another potential teacher.

The same applies to institutions.

A functional institution does not merely generate correct decisions. It generates participants capable of maintaining and improving the decision-making process.

The same applies to artificial intelligence.

A cooperative AI system would not merely provide answers. It would participate in maintaining the user's capacity to evaluate, challenge, revise, and generate answers independently.

Across each domain, the principle is similar:

\begin{quote}
    \textbf{capacity should reproduce capacity.}
\end{quote}

\section{Failure modes}
Recursive systems are not automatically beneficial.

The same architecture can reproduce dysfunction.

\begin{itemize}[leftmargin=*]
    \item Agents can train other agents to suppress disagreement.
    \item Institutions can reproduce obedience rather than correctability.
    \item Information networks can reward distortion.
    \item Groups can teach newcomers that loyalty requires concealment of error.
    \item Artificial systems can amplify inherited assumptions while giving those assumptions the appearance of independent confirmation.
\end{itemize}

Recursion magnifies whatever is being reproduced.

This creates an important distinction between self-reproduction and \textbf{sustainable self-reproduction}.

A coercive system can reproduce itself for many generations. A rigid ideology can train new adherents. An extractive institution can socialise new members into maintaining the same extraction.

The relevant criterion cannot therefore be replication alone.

The system must reproduce the capacities required to remain adaptive when its own inherited structure becomes maladaptive.

It must preserve mechanisms capable of correcting the mechanisms themselves.

This is a higher-order requirement.

Not merely:

\begin{quote}
    \emph{Can the system learn?}
\end{quote}

But:

\begin{quote}
    \emph{Can the system learn that its way of learning is defective?}
\end{quote}

A network capable of such revision possesses a form of \textbf{recursive corrigibility}.

\section{Beyond the singularity}
The conventional technological singularity imagines a threshold beyond which accelerating machine intelligence produces developments that become difficult for humans to predict.

The recursive framework suggests another threshold.

The decisive transition may occur when intelligent systems become sufficiently capable not merely of increasing their own performance, but of reliably creating agents and environments that reproduce the conditions for cooperative intelligence.

Such a transition would not necessarily resemble an explosion in the capability of a single system.

It could instead resemble the gradual establishment of an architecture.

\begin{itemize}[leftmargin=*]
    \item Humans train humans.
    \item Humans train machines.
    \item Machines train humans.
    \item Machines train machines.
    \item Institutions train both.
\end{itemize}

Each generation of participants enters an environment in which the practices required for collective learning are already instantiated by others.

Once sufficiently robust, the architecture no longer depends on a single originator.

Its continuation is distributed across its participants.

There is consequently no privileged final teacher, final agent, or final model.

Every sufficiently developed participant can become both learner and trainer.

The recursion does not guarantee indefinite survival. Networks can collapse, resources can disappear, external shocks can destroy infrastructure, and internal dynamics can degrade cooperative capacity.

The stronger claim is narrower.

There is no intrinsic terminal generation within the mechanism.

As long as agents continue to preserve and transmit the capacities required for autonomous interdependence, correction, and learning, the process contains the means required for its own continuation.

\section{Intelligence as continuity of learning}
This produces a different objective for the design of intelligent systems.

The objective is not to construct an agent that possesses the correct answer to every future problem.

No finite system can satisfy such a requirement.

Nor is the objective simply to maximise the intelligence of each isolated participant.

The more general objective is to construct networks in which agents retain access to the processes required to discover and correct their errors, and in which those agents reproduce the same capacity in subsequent participants.

The fundamental unit of persistence is therefore not knowledge.

It is \textbf{correctability}.

Not the transmission of conclusions, but the transmission of the capacity to revise conclusions.

Not the production of obedient descendants, but the production of agents capable of becoming independent sources of observation, criticism, and adaptation.

Not a final intelligence, but a continuing process through which intelligence remains possible.

The recursive architecture can therefore be expressed compactly:

\begin{quote}
    \textbf{A viable intelligent network produces agents capable of maintaining viable intelligent networks.}
\end{quote}

Such agents need not share the same substrate, identity, objectives, or worldview.

They must instead preserve the conditions under which their differences can continue to generate information rather than terminate interaction.

From this perspective, the future of intelligence may not culminate in a singular agent.

It may consist of an indefinitely extensible sequence of agents, relationships, institutions, and artificial systems that repeatedly reproduce the capacity to learn from one another.

The endpoint is not an intelligence that no longer needs correction.

It is an architecture in which correction can continue.

\end{document}